\documentclass[11pt,a4paper]{article}
\usepackage[utf8]{inputenc}
\usepackage[T1]{fontenc}
\usepackage[spanish,es-nodecimaldot]{babel}
\usepackage[margin=2.5cm]{geometry}
\usepackage{amsmath}
\usepackage{booktabs}
\usepackage{graphicx}
\usepackage{float}
\usepackage{tikz}
\usetikzlibrary{arrows.meta,babel}
\usepackage{microtype}
\usepackage{hyperref}
\hypersetup{hidelinks}
\setlength{\parindent}{0pt}
\setlength{\parskip}{0.6em}
\setlength{\emergencystretch}{2em}

\title{Práctica 3: Introducción a IRAF\\
\large Informe de avance}
\author{Astronomía Práctica II}
\date{}

\begin{document}
\maketitle

\section*{Objetivo y alcance del avance}
Explorar el manejo de imágenes astronómicas en IRAF y determinar la
traslación necesaria para alinear dos imágenes de un mismo campo estelar,
tomando \texttt{im010} como referencia. Este avance organiza el procedimiento
registrado y las mediciones de tres estrellas, siguiendo las fases 1 y 2
de la guía~\cite{guia}. Se completan las explicaciones del requisito de
entrega 1 y los cálculos del requisito 2. La ejecución del desplazamiento
calculado y su verificación mediante capturas de DS9 quedan pendientes de
documentar.

\section{Preparación, lectura de archivos y exploración de parámetros}

\subsection{Inicialización de la sesión}
Se inició IRAF desde la terminal y se abrió DS9 para visualizar las imágenes.
El símbolo \texttt{\$} identifica la terminal del sistema y \texttt{cl>}
identifica el intérprete de comandos de IRAF:
\begin{verbatim}
$ iraf
cl> !ds9 &
\end{verbatim}

\subsection{Preparación del directorio y lectura con \texttt{rfits}}
La siguiente transcripción organiza los comandos de los apuntes.
Las rutas suponen que se parte del directorio personal y que
\texttt{Documents} y \texttt{Downloads} son carpetas hermanas. Si un
directorio ya existe, se omite su creación. Desde
\texttt{Documents/ejercicio/intro} se requieren tres niveles de retorno
para acceder a \texttt{Downloads}; las rutas deben ajustarse si la
ubicación real es distinta.
\begin{verbatim}
cl> cd Documents
cl> mkdir ejercicio
cl> mkdir ejercicio/intro
cl> cd ejercicio/intro
cl> !mv ../../../Downloads/fintro0001 .
cl> !mv ../../../Downloads/fintro0002 .
cl> unlearn rfits
cl> rfits fintro* "" junk old+
cl> dir
\end{verbatim}

Aunque la guía denomina esta etapa ``descompresión'', \texttt{rfits}
lee datos FITS y los convierte en imágenes utilizables por IRAF; no es
un descompresor genérico. \texttt{unlearn rfits} restablece la configuración
de la tarea. Para archivos en disco, \verb|""| selecciona la unidad
primaria y \texttt{old+} activa \texttt{oldirafname}, que intenta recuperar
el nombre almacenado en la cabecera. Por ello, \texttt{junk} no determina
necesariamente el nombre final~\cite{iraf}. Se debe comprobar con
\texttt{dir} la presencia de las imágenes \texttt{im010} e \texttt{im011}
indicadas en la guía.

La figura~\ref{fig:rfits} desglosa el comando de lectura
\texttt{rfits} utilizado en este paso.

\begin{figure}[H]
\centering
\begin{tikzpicture}[
    font=\small,
    token/.style={draw=teal!65!black,fill=teal!7,rounded corners=3pt,
        minimum width=2.2cm,minimum height=0.75cm,font=\ttfamily},
    explanation/.style={anchor=west,text width=9.4cm,align=left},
    connection/.style={-{Stealth[length=2mm]},gray!70!black}]
\node[draw=gray!40,fill=gray!6,rounded corners=3pt,
    minimum width=13cm,minimum height=0.8cm] at (5.7,0)
    {\texttt{\detokenize{cl> rfits fintro* "" junk old+}}};

\node[token] (task) at (0.5,-1.2) {rfits};
\node[explanation] (tasktext) at (2.5,-1.2)
    {\textbf{Tarea.} Lee datos de imagen FITS y los convierte
    en imágenes utilizables por IRAF.};
\draw[connection] (task.east) -- (tasktext.west);

\node[token] (source) at (0.5,-2.5) {fintro*};
\node[explanation] (sourcetext) at (2.5,-2.5)
    {\textbf{Entrada: \texttt{fits\_file}.} Selecciona los archivos
    cuyo nombre comienza por \texttt{fintro}; por ejemplo,
    \texttt{fintro0001} y \texttt{fintro0002}.};
\draw[connection] (source.east) -- (sourcetext.west);

\node[token] (selection) at (0.5,-3.8) {\detokenize{""}};
\node[explanation] (selectiontext) at (2.5,-3.8)
    {\textbf{Selección: \texttt{file\_list}.} Para estos archivos
    en disco, la cadena vacía indica leer solo la unidad primaria
    (índice 0), no todas las extensiones.};
\draw[connection] (selection.east) -- (selectiontext.west);

\node[token] (output) at (0.5,-5.1) {junk};
\node[explanation] (outputtext) at (2.5,-5.1)
    {\textbf{Salida: \texttt{iraf\_file}.} Nombre raíz de salida.
    Si se usa esta raíz para varios archivos, se añaden números
    de secuencia. No es un archivo de entrada.};
\draw[connection] (output.east) -- (outputtext.west);

\node[token] (names) at (0.5,-6.4) {old+};
\node[explanation] (namestext) at (2.5,-6.4)
    {\textbf{Nombres originales.} Abrevia \texttt{oldirafname=yes}.
    Intenta usar el nombre guardado en la palabra clave
    \texttt{IRAFNAME} de la cabecera, en lugar de la raíz \texttt{junk}.};
\draw[connection] (names.east) -- (namestext.west);
\end{tikzpicture}
\caption{Desglose del comando \texttt{rfits} para los archivos en disco
 de la práctica. \texttt{cl>} es el indicador del intérprete, no un
 argumento. La recuperación de los nombres originales depende de la
 información disponible en las cabeceras~\cite{iraf}.}
\label{fig:rfits}
\end{figure}


\subsection{Exploración y modificación de parámetros}
Se comparó el resumen de una imagen con su cabecera completa y se exploró
la configuración de \texttt{imheader}, abreviada como \texttt{imhead}:
\begin{verbatim}
cl> lpar imheader
cl> imhead im010
cl> imhead im010 longheader=yes | page
cl> epar imhead
\end{verbatim}
En el editor se asigna \texttt{yes} a \texttt{longheader}, se confirma
la entrada y se sale guardando con \texttt{:q}. Luego se verifica la
configuración con:
\begin{verbatim}
cl> lpar imhead
\end{verbatim}
La edición con \texttt{epar} permite conservar el cambio para ejecuciones
posteriores~\cite{iraf}.

\subsection{Requisito de entrega 1: fundamentos}

\textbf{Pendiente:} reformular esta respuesta de manera más breve y concisa.
Por ahora, se conserva el texto sin modificaciones.

\subsubsection*{Parámetros Query y Hidden}
Los parámetros \emph{Query} se consultan al usuario cuando la tarea
necesita su valor y no se han especificado en la línea de comandos.
En \texttt{lpar} aparecen sin paréntesis. Los parámetros \emph{Hidden}
aparecen entre paréntesis y normalmente emplean su valor configurado sin
preguntar; pueden indicarse explícitamente como \texttt{nombre=valor}.
Una excepción ocurre si el valor oculto es indefinido o inválido.
Por tanto, ``obligatorio'' y ``opcional'' describen aquí la interacción
con el usuario, no la importancia física del parámetro~\cite{iraf}.

En este caso, \texttt{im010} identifica los datos que se examinan,
mientras que \texttt{longheader=yes} controla la información mostrada:
permite inspeccionar la cabecera completa sin modificar los píxeles.
Especificar un parámetro oculto en una llamada no suele guardar su valor
para la siguiente; editarlo con \texttt{epar} sí permite persistir el
cambio~\cite{iraf}.

\subsubsection*{Justificación física del formato FITS}
En astrofísica se necesita conservar los valores numéricos de los píxeles,
no solamente su apariencia. FITS organiza los datos en unidades de
cabecera y datos. Las cabeceras contienen registros ASCII legibles con
palabras clave, valores y comentarios; permiten documentar dimensiones,
tipo numérico, condiciones de observación e historial de procesamiento.
Así se conserva el contexto necesario para interpretar y reproducir una
medición científica~\cite{fitsprimer}.

El campo \texttt{BITPIX} especifica la representación de los datos:
FITS admite enteros de 8, 16, 32 y 64 bits y números de coma flotante
de 32 o 64 bits, estos últimos identificados por \texttt{BITPIX=-32}
y \texttt{BITPIX=-64}. Esta flexibilidad permite conservar un rango
numérico amplio y resultados fraccionarios o negativos de la reducción,
si se elige el tipo adecuado; guardar en FITS no recupera información
que ya se haya perdido~\cite{fitsstandard}.

El JPEG convencional utiliza compresión con pérdidas: la cuantización
de su representación transformada altera los valores reconstruidos.
Por ello, una imagen visualmente similar no garantiza las mismas sumas
de intensidad ni los mismos perfiles estelares. Esto se refiere al JPEG
de uso habitual, no a todos los modos de la familia JPEG~\cite{jpeg}.
En una representación de 8 bits por canal solo hay $2^8=256$ niveles,
frente a $2^{16}=65536$ en una representación entera de 16 bits.

PNG, en cambio, \textbf{sí comprime sin pérdidas} y admite hasta 16 bits
por muestra. También puede almacenar metadatos textuales. Su limitación
no es una pérdida inherente por compresión, sino que no ofrece la misma
representación de datos científicos en coma flotante ni las convenciones
astronómicas de FITS. Si se exporta una visualización previamente
reescalada, recortada o reducida a 8 bits, PNG conserva esa visualización,
no los datos científicos originales~\cite{png,fitsstandard}.
En consecuencia, JPG y PNG son útiles para ilustrar el informe, mientras
que los cálculos deben realizarse sobre las imágenes científicas.

\section{Medición de centroides y alineamiento estelar}

\subsection{Visualización y parpadeo en DS9}
Se cargaron ambas imágenes en marcos distintos:
\begin{verbatim}
cl> display im010 1
cl> display im011 2
\end{verbatim}

La opción \emph{Match LUTs} se refiere a igualar las tablas de
transferencia de color, es decir, la correspondencia entre los niveles
de visualización y los colores o tonos de gris. Está relacionada con
la barra de color, pero no debe confundirse con igualar únicamente los
límites mínimo y máximo de intensidad: DS9 distingue el emparejamiento
de la barra de color, de la escala y de sus límites~\cite{ds9}.
Estas opciones afectan la visualización, no corrigen las posiciones
estelares en los datos.

El modo \emph{Blink} alterna los marcos y facilita identificar el
desplazamiento aparente de las mismas estrellas entre ambas imágenes,
como propone la guía~\cite{guia}. El parpadeo se desactiva antes de
realizar las mediciones interactivas.

\subsection{Medición con \texttt{imexamine}}
Se inició la herramienta mediante:
\begin{verbatim}
cl> unlearn imexamine
cl> imexamine
\end{verbatim}
Para obtener las posiciones se sitúa el cursor sobre cada estrella de
referencia y se utiliza la tecla \texttt{a}; se repite el procedimiento
para los mismos objetos en la otra imagen. La tecla \texttt{d} permite
solicitar la carga de una imagen, cuyo nombre y marco deben indicarse;
no selecciona por sí sola \texttt{im011}. Se sale con
\texttt{q}~\cite{iraf}.

El centroide es la posición estimada del centro de la distribución
estelar, mientras que la PSF (\emph{Point Spread Function}) describe
su perfil espacial. No son la misma magnitud. La tarea dispone de
centrado y análisis de perfiles, y el modelo de ajuste depende de sus
parámetros; sin el registro de esa configuración no se afirma que estas
mediciones procedan de un ajuste gaussiano bidimensional~\cite{iraf}.

\subsection{Requisito de entrega 2: tabla y desplazamiento promedio}
Se organizan los datos anotados y la corrección indicada para la estrella 1,
interpretando las cuatro columnas como $X_{10}$, $Y_{10}$, $X_{11}$ y $Y_{11}$, en píxeles.
Para que los deltas sean directamente la corrección aplicada a
\texttt{im011}, se adopta la convención de la guía:
\begin{equation}
\Delta x_i=X_{10,i}-X_{11,i},\qquad
\Delta y_i=Y_{10,i}-Y_{11,i}.
\label{eq:deltas}
\end{equation}

\begin{table}[htbp]
\centering
\small
\setlength{\tabcolsep}{5pt}
\caption{Alineamiento estelar con la diferencia vertical corregida de la estrella 1.
Todas las coordenadas y diferencias están expresadas en píxeles.}
\label{tab:alineamiento}
\begin{tabular}{crrrrrr}
\toprule
Estrella & $X_{10}$ & $Y_{10}$ & $X_{11}$ & $Y_{11}$
& $\Delta x$ & $\Delta y$ \\
\midrule
1 & 289.986 & 83.967  & 289.986 & 85.547  & 0.000  & $-1.580$ \\
2 & 281.306 & 189.292 & 281.884 & 191.028 & $-0.578$ & $-1.736$ \\
3 & 331.347 & 134.894 & 332.320 & 136.630 & $-0.973$ & $-1.736$ \\
\midrule
\multicolumn{5}{l}{Promedio de las diferencias}
& $-0.5170$ & $-1.6840$ \\
\bottomrule
\end{tabular}
\end{table}

\subsection{Comando de corrección con \texttt{imshift}}
IRAF define la transformación como
$x_{\mathrm{salida}}=x_{\mathrm{entrada}}+\texttt{xshift}$ y
$y_{\mathrm{salida}}=y_{\mathrm{entrada}}+\texttt{yshift}$~\cite{iraf}.
Por tanto, el comando completado a partir de los datos registrados
en la tabla~\ref{tab:alineamiento} es:
\begin{verbatim}
cl> unlearn imshift
cl> imshift im011 s011 -0.5170 -1.6840
\end{verbatim}
\textbf{Estado:} comando calculado, pendiente de confirmar como ejecutado.
Los signos negativos desplazan las estrellas hacia coordenadas menores
para compensar el desfase medio de \texttt{im011}. Al ser un desplazamiento
fraccionario, se requiere interpolación; por defecto, tras
\texttt{unlearn}, \texttt{imshift} utiliza interpolación
bilineal~\cite{iraf}.

La figura~\ref{fig:imshift} ilustra la traslación aplicada a la estrella 1
con los desplazamientos promedio de la tabla~\ref{tab:alineamiento}.

\begin{figure}[H]
\centering
\begin{tikzpicture}[x=3cm,y=3cm,font=\small,
    >=Stealth,
    initial/.style={blue!70!black},
    shifted/.style={orange!85!black},
    correction/.style={teal!75!black}]
\coordinate (reference) at (0,0);
\coordinate (initial) at (0,1.580);
\coordinate (shifted) at (-0.517,-0.104);
\coordinate (corner) at (-0.517,1.580);

\draw[step=0.5,gray!20,very thin] (-1,-0.5) grid (0.75,2);
\draw[->,gray!70!black] (-1,-0.5) -- (0.9,-0.5);
\draw[->,gray!70!black] (-1,-0.5) -- (-1,2.12);
\foreach \tick in {-1,-0.5,0,0.5}{
    \draw[gray!70!black] (\tick,-0.5) -- (\tick,-0.54)
        node[below,font=\footnotesize] {$\tick$};
}
\foreach \tick in {0,0.5,1,1.5,2}{
    \draw[gray!70!black] (-1,\tick) -- (-1.04,\tick)
        node[left,font=\footnotesize] {$\tick$};
}
\node[below] at (-0.08,-0.74) {$x-X_{10,1}$ (píxeles)};
\node[rotate=90] at (-1.42,0.75) {$y-Y_{10,1}$ (píxeles)};

\draw[correction,dashed,->] (initial) -- (corner)
    node[midway,above=5pt] {$\overline{\Delta x}=-0.517$};
\draw[correction,dashed,->] (corner) -- (shifted)
    node[midway,left=5pt,rotate=90,anchor=south]
    {$\overline{\Delta y}=-1.684$};
\draw[correction,very thick,->,shorten <=3pt,shorten >=4pt]
    (initial) -- (shifted);
\node[correction,anchor=west,align=left] at (0.08,0.85)
    {Traslación\\\texttt{imshift}};

\draw[gray!70!black,densely dotted,thick] (shifted) -- (reference);
\fill[initial] (initial) circle (2.6pt);
\node[initial,right=6pt] at (initial) {\texttt{im011}};
\fill[shifted] (shifted) circle (2.6pt);
\node[shifted,below left=3pt] at (shifted) {\texttt{s011}};
\draw[black,thick] (-0.055,0) -- (0.055,0);
\draw[black,thick] (0,-0.055) -- (0,0.055);
\node[above right=3pt] at (reference) {\texttt{im010}};
\end{tikzpicture}

\smallskip
{\small
\begin{tabular}{lrr}
\toprule
Posición de la estrella 1 & $X$ (píxeles) & $Y$ (píxeles) \\
\midrule
Referencia: \texttt{im010} & 289.986 & 83.967 \\
Entrada: \texttt{im011} & 289.986 & 85.547 \\
Salida predicha: \texttt{s011} & 289.469 & 83.863 \\
\bottomrule
\end{tabular}
}
\caption{Esquema geométrico de la corrección con \texttt{imshift},
aplicada a la estrella 1. El origen de las coordenadas relativas es su
posición en \texttt{im010}.}
\label{fig:imshift}
\end{figure}

\subsection{Comprobación del resultado y capturas pendientes}
Después de ejecutar la corrección, se deben cargar la imagen de referencia
y la imagen desplazada (con blink):

% \begin{figure}[htbp]
% \centering
% \includegraphics[width=0.9\linewidth]{Blink1.jpg}
% \caption{Primer estado del parpadeo en DS9: imagen de referencia
% \texttt{im010}, con ejes en píxeles. Comparar las posiciones de las
% estrellas con el segundo estado, conservando el encuadre y la escala.}
% \label{fig:blink1}
% \end{figure}

% \begin{figure}[htbp]
% \centering
% \includegraphics[width=0.9\linewidth]{Blink2.jpg}
% \caption{Segundo estado del parpadeo en DS9: imagen desplazada
% \texttt{s011}, con ejes en píxeles. La comparación con \texttt{im010}
% permite evaluar visualmente el alineamiento de las estrellas.}
% \label{fig:blink2}
% \end{figure}


\bibliographystyle{unsrt}
\bibliography{bibliography}

\end{document}